% CHƯƠNG 4: sáu thành phần cố định cho mỗi nội dung; hình/bảng chèn tự động; ô TODO do nhóm điền.
\clearpage
\section[Kiểm chứng chuỗi dừng và mô hình ARMA]{KIỂM CHỨNG CHUỖI DỪNG VÀ MÔ HÌNH ARMA}

\TODO{Đoạn dẫn: các bộ dữ liệu áp dụng (Hồ Huron, Linh miêu, DAX suất sinh lợi); cách đọc một nội dung; câu nối sang Chương 5.}

\subsection{Chuỗi dừng}

\noidung{Tính dừng và nhiễu trắng}
\tp{Nội dung lý thuyết cần kiểm chứng} \TODO{một câu mệnh đề cần kiểm chứng; dẫn mục 2.1.1(a)}
\tp{Cách triển khai} Dữ liệu thực. \textit{Bộ áp dụng:} Hồ Huron, Linh miêu, DAX (suất sinh lợi). \TODO{quy trình cụ thể, tham số riêng của mục này}
\tp{Mã lệnh Python} \TODO{đoạn mã lõi 10--20 dòng (chỉ ký tự ASCII và chữ Việt, không dùng ký hiệu Hy Lạp trong mã); mã đầy đủ ở notebook}
\tp{Kết quả tính toán}
\input{\tabdir/tab_b11_01_stationarity.tex}
\figx{fig_b11_01_series_rolling}{Chuỗi, trung bình trượt (trái) và phương sai trượt (phải) của HU, LY, DR}
\TODO{chọn bảng/hình giữ lại, bỏ bớt phần phụ; viết câu dẫn}
\tp{Phân tích và diễn giải}
\bodl{(i) Hồ Huron.} \TODO{diễn giải gắn với 2.1.1(a)}

\bodl{(ii) Linh miêu.} \TODO{diễn giải gắn với 2.1.1(a)}

\bodl{(iii) DAX (suất sinh lợi).} \TODO{diễn giải gắn với 2.1.1(a)}

\tp{Nhận xét} \TODO{xếp loại (phù hợp / phù hợp một phần / không phù hợp); nếu sai lệch thì nêu nguyên nhân}

\noidung{Tính chất của tự hiệp phương sai}
\tp{Nội dung lý thuyết cần kiểm chứng} \TODO{một câu mệnh đề cần kiểm chứng; dẫn mục 2.1.1(b)}
\tp{Cách triển khai} Dữ liệu thực. \textit{Bộ áp dụng:} Hồ Huron, Linh miêu, DAX (suất sinh lợi). \TODO{quy trình cụ thể, tham số riêng của mục này}
\tp{Mã lệnh Python} \TODO{đoạn mã lõi 10--20 dòng (chỉ ký tự ASCII và chữ Việt, không dùng ký hiệu Hy Lạp trong mã); mã đầy đủ ở notebook}
\tp{Kết quả tính toán}
\input{\tabdir/tab_b11_02_gamma_properties.tex}
\figx{fig_b11_02_acf_denominators}{$\hat\rho(h)$ với mẫu số $n$ và $n-h$ ở mọi độ trễ $h=1,\dots,n-1$}
\TODO{chọn bảng/hình giữ lại, bỏ bớt phần phụ; viết câu dẫn}
\tp{Phân tích và diễn giải}
\bodl{(i) Hồ Huron.} \TODO{diễn giải gắn với 2.1.1(b)}

\bodl{(ii) Linh miêu.} \TODO{diễn giải gắn với 2.1.1(b)}

\bodl{(iii) DAX (suất sinh lợi).} \TODO{diễn giải gắn với 2.1.1(b)}

\tp{Nhận xét} \TODO{xếp loại (phù hợp / phù hợp một phần / không phù hợp); nếu sai lệch thì nêu nguyên nhân}

\noidung{Trung bình mẫu của chuỗi phụ thuộc}
\tp{Nội dung lý thuyết cần kiểm chứng} \TODO{một câu mệnh đề cần kiểm chứng; dẫn mục 2.1.1(c)}
\tp{Cách triển khai} Cả hai. \textit{Bộ áp dụng:} Hồ Huron. \textit{Không áp dụng cho:} Linh miêu, DAX (suất sinh lợi) \TODO{lý do ngắn}. \TODO{quy trình cụ thể, tham số riêng của mục này}
\tp{Mã lệnh Python} \TODO{đoạn mã lõi 10--20 dòng (chỉ ký tự ASCII và chữ Việt, không dùng ký hiệu Hy Lạp trong mã); mã đầy đủ ở notebook}
\tp{Kết quả tính toán}
\input{\tabdir/tab_b11_03_coverage.tex}
\input{\tabdir/tab_b11_03_mean_se.tex}
\input{\tabdir/tab_b11_03_slope_se.tex}
\TODO{chọn bảng/hình giữ lại, bỏ bớt phần phụ; viết câu dẫn}
\tp{Phân tích và diễn giải}
\bodl{(i) Hồ Huron.} \TODO{diễn giải gắn với 2.1.1(c)}

\tp{Nhận xét} \TODO{xếp loại (phù hợp / phù hợp một phần / không phù hợp); nếu sai lệch thì nêu nguyên nhân}

\noidung{Quá trình tuyến tính}
\tp{Nội dung lý thuyết cần kiểm chứng} \TODO{một câu mệnh đề cần kiểm chứng; dẫn mục 2.1.1(d)}
\tp{Cách triển khai} Dữ liệu thực. \textit{Bộ áp dụng:} Hồ Huron, Linh miêu. \textit{Không áp dụng cho:} DAX (suất sinh lợi) \TODO{lý do ngắn}. \TODO{quy trình cụ thể, tham số riêng của mục này}
\tp{Mã lệnh Python} \TODO{đoạn mã lõi 10--20 dòng (chỉ ký tự ASCII và chữ Việt, không dùng ký hiệu Hy Lạp trong mã); mã đầy đủ ở notebook}
\tp{Kết quả tính toán}
\input{\tabdir/tab_b11_05_psi_acf.tex}
\figx{fig_b11_05_acf_theory}{ACF mẫu (que) và ACF lý thuyết của AR(2) khớp bằng ML (đường)}
\TODO{chọn bảng/hình giữ lại, bỏ bớt phần phụ; viết câu dẫn}
\tp{Phân tích và diễn giải}
\bodl{(i) Hồ Huron.} \TODO{diễn giải gắn với 2.1.1(d)}

\bodl{(ii) Linh miêu.} \TODO{diễn giải gắn với 2.1.1(d)}

\tp{Nhận xét} \TODO{xếp loại (phù hợp / phù hợp một phần / không phù hợp); nếu sai lệch thì nêu nguyên nhân}

\subsection{Mô hình ARMA}

\noidung{Mô hình AR(1)}
\tp{Nội dung lý thuyết cần kiểm chứng} \TODO{một câu mệnh đề cần kiểm chứng; dẫn mục 2.1.2(a)}
\tp{Cách triển khai} Dữ liệu thực. \textit{Bộ áp dụng:} Hồ Huron. \textit{Không áp dụng cho:} Linh miêu, DAX (suất sinh lợi) \TODO{lý do ngắn}. \TODO{quy trình cụ thể, tham số riêng của mục này}
\tp{Mã lệnh Python} \TODO{đoạn mã lõi 10--20 dòng (chỉ ký tự ASCII và chữ Việt, không dùng ký hiệu Hy Lạp trong mã); mã đầy đủ ở notebook}
\tp{Kết quả tính toán}
\input{\tabdir/tab_b11_06_ar1_acf.tex}
\input{\tabdir/tab_b11_06_ar1_forecast_mse.tex}
\input{\tabdir/tab_b11_06_ar1_simulation.tex}
\TODO{chọn bảng/hình giữ lại, bỏ bớt phần phụ; viết câu dẫn}
\tp{Phân tích và diễn giải}
\bodl{(i) Hồ Huron.} \TODO{diễn giải gắn với 2.1.2(a)}

\tp{Nhận xét} \TODO{xếp loại (phù hợp / phù hợp một phần / không phù hợp); nếu sai lệch thì nêu nguyên nhân}

\noidung{Mô hình AR(p) và hệ Yule–Walker}
\tp{Nội dung lý thuyết cần kiểm chứng} \TODO{một câu mệnh đề cần kiểm chứng; dẫn mục 2.1.2(b)}
\tp{Cách triển khai} Dữ liệu thực. \textit{Bộ áp dụng:} Linh miêu, Hồ Huron. \textit{Không áp dụng cho:} DAX (suất sinh lợi) \TODO{lý do ngắn}. \TODO{quy trình cụ thể, tham số riêng của mục này}
\tp{Mã lệnh Python} \TODO{đoạn mã lõi 10--20 dòng (chỉ ký tự ASCII và chữ Việt, không dùng ký hiệu Hy Lạp trong mã); mã đầy đủ ở notebook}
\tp{Kết quả tính toán}
\input{\tabdir/tab_b11_07_yule_walker.tex}
\TODO{chọn bảng/hình giữ lại, bỏ bớt phần phụ; viết câu dẫn}
\tp{Phân tích và diễn giải}
\bodl{(i) Linh miêu.} \TODO{diễn giải gắn với 2.1.2(b)}

\bodl{(ii) Hồ Huron.} \TODO{diễn giải gắn với 2.1.2(b)}

\tp{Nhận xét} \TODO{xếp loại (phù hợp / phù hợp một phần / không phù hợp); nếu sai lệch thì nêu nguyên nhân}

\noidung{Mô hình AR(2)}
\tp{Nội dung lý thuyết cần kiểm chứng} \TODO{một câu mệnh đề cần kiểm chứng; dẫn mục 2.1.2(c)}
\tp{Cách triển khai} Dữ liệu thực. \textit{Bộ áp dụng:} Linh miêu. \textit{Không áp dụng cho:} Hồ Huron, DAX (suất sinh lợi) \TODO{lý do ngắn}. \TODO{quy trình cụ thể, tham số riêng của mục này}
\tp{Mã lệnh Python} \TODO{đoạn mã lõi 10--20 dòng (chỉ ký tự ASCII và chữ Việt, không dùng ký hiệu Hy Lạp trong mã); mã đầy đủ ở notebook}
\tp{Kết quả tính toán}
\input{\tabdir/tab_b11_08_ar2_triangle.tex}
\figx{fig_b11_08_ar2_triangle}{Tam giác nhân quả AR(2) với điểm HU, LY (trái) và ACF mẫu của LY so với AR(2)-YW (phải)}
\TODO{chọn bảng/hình giữ lại, bỏ bớt phần phụ; viết câu dẫn}
\tp{Phân tích và diễn giải}
\bodl{(i) Linh miêu.} \TODO{diễn giải gắn với 2.1.2(c)}

\tp{Nhận xét} \TODO{xếp loại (phù hợp / phù hợp một phần / không phù hợp); nếu sai lệch thì nêu nguyên nhân}

\noidung{Mô hình MA(q)}
\tp{Nội dung lý thuyết cần kiểm chứng} \TODO{một câu mệnh đề cần kiểm chứng; dẫn mục 2.1.2(d)}
\tp{Cách triển khai} Cả hai. \textit{Bộ áp dụng:} DAX (suất sinh lợi), Hồ Huron. \textit{Không áp dụng cho:} Linh miêu \TODO{lý do ngắn}. \TODO{quy trình cụ thể, tham số riêng của mục này}
\tp{Mã lệnh Python} \TODO{đoạn mã lõi 10--20 dòng (chỉ ký tự ASCII và chữ Việt, không dùng ký hiệu Hy Lạp trong mã); mã đầy đủ ở notebook}
\tp{Kết quả tính toán}
\input{\tabdir/tab_b11_09_ma1.tex}
\input{\tabdir/tab_b11_09_ma2_simulation.tex}
\figx{fig_b11_09_ma2_acf}{MA(2) mô phỏng: trung bình ACF mẫu so với ACF lý thuyết cắt cụt sau độ trễ 2}
\TODO{chọn bảng/hình giữ lại, bỏ bớt phần phụ; viết câu dẫn}
\tp{Phân tích và diễn giải}
\bodl{(i) DAX (suất sinh lợi).} \TODO{diễn giải gắn với 2.1.2(d)}

\bodl{(ii) Hồ Huron.} \TODO{diễn giải gắn với 2.1.2(d)}

\tp{Nhận xét} \TODO{xếp loại (phù hợp / phù hợp một phần / không phù hợp); nếu sai lệch thì nêu nguyên nhân}

\noidung{Mô hình ARMA(p, q)}
\tp{Nội dung lý thuyết cần kiểm chứng} \TODO{một câu mệnh đề cần kiểm chứng; dẫn mục 2.1.2(e)}
\tp{Cách triển khai} Dữ liệu thực. \textit{Bộ áp dụng:} Hồ Huron, Linh miêu. \textit{Không áp dụng cho:} DAX (suất sinh lợi) \TODO{lý do ngắn}. \TODO{quy trình cụ thể, tham số riêng của mục này}
\tp{Mã lệnh Python} \TODO{đoạn mã lõi 10--20 dòng (chỉ ký tự ASCII và chữ Việt, không dùng ký hiệu Hy Lạp trong mã); mã đầy đủ ở notebook}
\tp{Kết quả tính toán}
\input{\tabdir/tab_b11_10_arma11.tex}
\input{\tabdir/tab_b11_10_cancellation.tex}
\input{\tabdir/tab_b11_10_overparameterized.tex}
\TODO{chọn bảng/hình giữ lại, bỏ bớt phần phụ; viết câu dẫn}
\tp{Phân tích và diễn giải}
\bodl{(i) Hồ Huron.} \TODO{diễn giải gắn với 2.1.2(e)}

\bodl{(ii) Linh miêu.} \TODO{diễn giải gắn với 2.1.2(e)}

\tp{Nhận xét} \TODO{xếp loại (phù hợp / phù hợp một phần / không phù hợp); nếu sai lệch thì nêu nguyên nhân}

\subsection{Xây dựng mô hình ARMA}

\noidung{Tự tương quan và kiểm định Ljung–Box}
\tp{Nội dung lý thuyết cần kiểm chứng} \TODO{một câu mệnh đề cần kiểm chứng; dẫn mục 2.1.3(a)}
\tp{Cách triển khai} Cả hai. \textit{Bộ áp dụng:} Hồ Huron, Linh miêu, DAX (suất sinh lợi). \TODO{quy trình cụ thể, tham số riêng của mục này}
\tp{Mã lệnh Python} \TODO{đoạn mã lõi 10--20 dòng (chỉ ký tự ASCII và chữ Việt, không dùng ký hiệu Hy Lạp trong mã); mã đầy đủ ở notebook}
\tp{Kết quả tính toán}
\input{\tabdir/tab_b11_04_acf_ljungbox.tex}
\input{\tabdir/tab_b11_04_simulation.tex}
\figx{fig_b11_04_acf_bands}{ACF mẫu và dải 95\% (đứt nét: iid $\pm1.96/\sqrt n$; chấm: Bartlett MA(1)); độ trễ vượt dải tô cam đỏ}
\TODO{chọn bảng/hình giữ lại, bỏ bớt phần phụ; viết câu dẫn}
\tp{Phân tích và diễn giải}
\bodl{(i) Hồ Huron.} \TODO{diễn giải gắn với 2.1.3(a)}

\bodl{(ii) Linh miêu.} \TODO{diễn giải gắn với 2.1.3(a)}

\bodl{(iii) DAX (suất sinh lợi).} \TODO{diễn giải gắn với 2.1.3(a)}

\tp{Nhận xét} \TODO{xếp loại (phù hợp / phù hợp một phần / không phù hợp); nếu sai lệch thì nêu nguyên nhân}

\noidung{Tự tương quan riêng phần}
\tp{Nội dung lý thuyết cần kiểm chứng} \TODO{một câu mệnh đề cần kiểm chứng; dẫn mục 2.1.3(b)}
\tp{Cách triển khai} Dữ liệu thực. \textit{Bộ áp dụng:} Hồ Huron, Linh miêu, DAX (suất sinh lợi). \TODO{quy trình cụ thể, tham số riêng của mục này}
\tp{Mã lệnh Python} \TODO{đoạn mã lõi 10--20 dòng (chỉ ký tự ASCII và chữ Việt, không dùng ký hiệu Hy Lạp trong mã); mã đầy đủ ở notebook}
\tp{Kết quả tính toán}
\input{\tabdir/tab_b11_11_pacf.tex}
\figx{fig_b11_11_pacf}{PACF mẫu (Durbin--Levinson) của HU đã khử xu thế, LY và DR}
\TODO{chọn bảng/hình giữ lại, bỏ bớt phần phụ; viết câu dẫn}
\tp{Phân tích và diễn giải}
\bodl{(i) Hồ Huron.} \TODO{diễn giải gắn với 2.1.3(b)}

\bodl{(ii) Linh miêu.} \TODO{diễn giải gắn với 2.1.3(b)}

\bodl{(iii) DAX (suất sinh lợi).} \TODO{diễn giải gắn với 2.1.3(b)}

\tp{Nhận xét} \TODO{xếp loại (phù hợp / phù hợp một phần / không phù hợp); nếu sai lệch thì nêu nguyên nhân}

\noidung{Ước lượng tham số}
\tp{Nội dung lý thuyết cần kiểm chứng} \TODO{một câu mệnh đề cần kiểm chứng; dẫn mục 2.1.3(c)}
\tp{Cách triển khai} Cả hai. \textit{Bộ áp dụng:} Linh miêu, Hồ Huron. \textit{Không áp dụng cho:} DAX (suất sinh lợi) \TODO{lý do ngắn}. \TODO{quy trình cụ thể, tham số riêng của mục này}
\tp{Mã lệnh Python} \TODO{đoạn mã lõi 10--20 dòng (chỉ ký tự ASCII và chữ Việt, không dùng ký hiệu Hy Lạp trong mã); mã đầy đủ ở notebook}
\tp{Kết quả tính toán}
\input{\tabdir/tab_b11_12_estimators.tex}
\input{\tabdir/tab_b11_12_simulation.tex}
\TODO{chọn bảng/hình giữ lại, bỏ bớt phần phụ; viết câu dẫn}
\tp{Phân tích và diễn giải}
\bodl{(i) Linh miêu.} \TODO{diễn giải gắn với 2.1.3(c)}

\bodl{(ii) Hồ Huron.} \TODO{diễn giải gắn với 2.1.3(c)}

\tp{Nhận xét} \TODO{xếp loại (phù hợp / phù hợp một phần / không phù hợp); nếu sai lệch thì nêu nguyên nhân}

\noidung{Lựa chọn bậc}
\tp{Nội dung lý thuyết cần kiểm chứng} \TODO{một câu mệnh đề cần kiểm chứng; dẫn mục 2.1.3(d)}
\tp{Cách triển khai} Dữ liệu thực. \textit{Bộ áp dụng:} Linh miêu, Hồ Huron. \textit{Không áp dụng cho:} DAX (suất sinh lợi) \TODO{lý do ngắn}. \TODO{quy trình cụ thể, tham số riêng của mục này}
\tp{Mã lệnh Python} \TODO{đoạn mã lõi 10--20 dòng (chỉ ký tự ASCII và chữ Việt, không dùng ký hiệu Hy Lạp trong mã); mã đầy đủ ở notebook}
\tp{Kết quả tính toán}
\input{\tabdir/tab_b11_13_dr_grid.tex}
\input{\tabdir/tab_b11_13_hu_models.tex}
\input{\tabdir/tab_b11_13_ly_ic.tex}
\input{\tabdir/tab_b11_13_ly_rolling.tex}
\figx{fig_b11_13_ic_ly}{LY: AIC và BIC (trừ giá trị nhỏ nhất) theo bậc AR; dấu sao là bậc được chọn}
\TODO{chọn bảng/hình giữ lại, bỏ bớt phần phụ; viết câu dẫn}
\tp{Phân tích và diễn giải}
\bodl{(i) Linh miêu.} \TODO{diễn giải gắn với 2.1.3(d)}

\bodl{(ii) Hồ Huron.} \TODO{diễn giải gắn với 2.1.3(d)}

\tp{Nhận xét} \TODO{xếp loại (phù hợp / phù hợp một phần / không phù hợp); nếu sai lệch thì nêu nguyên nhân}

\noidung{Chẩn đoán phần dư}
\tp{Nội dung lý thuyết cần kiểm chứng} \TODO{một câu mệnh đề cần kiểm chứng; dẫn mục 2.1.3(e)}
\tp{Cách triển khai} Dữ liệu thực. \textit{Bộ áp dụng:} Hồ Huron, Linh miêu, DAX (suất sinh lợi). \TODO{quy trình cụ thể, tham số riêng của mục này}
\tp{Mã lệnh Python} \TODO{đoạn mã lõi 10--20 dòng (chỉ ký tự ASCII và chữ Việt, không dùng ký hiệu Hy Lạp trong mã); mã đầy đủ ở notebook}
\tp{Kết quả tính toán}
\input{\tabdir/tab_b11_14_residual_diagnostics.tex}
\figx{fig_b11_14_residual_acf}{ACF của phần dư (trái) và của phần dư bình phương (phải) cho các mô hình được chọn}
\TODO{chọn bảng/hình giữ lại, bỏ bớt phần phụ; viết câu dẫn}
\tp{Phân tích và diễn giải}
\bodl{(i) Hồ Huron.} \TODO{diễn giải gắn với 2.1.3(e)}

\bodl{(ii) Linh miêu.} \TODO{diễn giải gắn với 2.1.3(e)}

\bodl{(iii) DAX (suất sinh lợi).} \TODO{diễn giải gắn với 2.1.3(e)}

\tp{Nhận xét} \TODO{xếp loại (phù hợp / phù hợp một phần / không phù hợp); nếu sai lệch thì nêu nguyên nhân}
